\documentclass[11pt]{article}
\usepackage[a4paper,margin=25mm]{geometry}
\usepackage{amsmath,amssymb,amsthm,mathtools,bm,slashed}
\usepackage{booktabs,longtable,array}
\usepackage{enumitem}
\usepackage{xcolor}
\usepackage{microtype}
\usepackage[hidelinks]{hyperref}

\setlength{\emergencystretch}{3em}
\newtheorem{theorem}{Theorem}[section]
\newtheorem{proposition}[theorem]{Proposition}
\newtheorem{lemma}[theorem]{Lemma}
\newtheorem{corollary}[theorem]{Corollary}
\theoremstyle{definition}
\newtheorem{definition}[theorem]{Definition}
\newtheorem{gate}[theorem]{Fail-closed gate}
\theoremstyle{remark}
\newtheorem{remark}[theorem]{Remark}

\newcommand{\ii}{\mathrm i}
\newcommand{\dd}{\mathrm d}
\newcommand{\CC}{\mathbb C}
\newcommand{\CP}{\mathbb {CP}}
\newcommand{\cO}{\mathcal O}
\newcommand{\cH}{\mathcal H}
\newcommand{\Tr}{\operatorname{Tr}}
\newcommand{\Sym}{\operatorname{Sym}}
\newcommand{\Span}{\operatorname{Span}}
\newcommand{\Ker}{\operatorname{Ker}}
\newcommand{\ind}{\operatorname{ind}}
\newcommand{\status}[1]{\textcolor{blue!60!black}{\textsf{#1}}}
\newcolumntype{L}[1]{>{\raggedright\arraybackslash}p{#1}}

\title{AP-E4: Tangent-Valued Dirac Operators on \(\mathbb{CP}^1\)\\
Chirality, Complete Partner Spectrum, Spectral Gap,\\
and the Canonical-Spin\texorpdfstring{\(^c\)}{c} versus Ordinary-Spin Gate}
\author{Route F research note}
\date{15 July 2026}

\begin{document}
\maketitle

\begin{abstract}
This note independently constructs and solves the AP-E4 operator problem.
Linearizing a \(\CP^1\) sigma-model field proves that its physical horizontal
variation is tangent-valued.  In affine charts the coefficient transforms by
\(\delta(1/w)=-w^{-2}\delta w\), hence
\(T^{1,0}\CP^1\simeq\cO(2)\) with Chern number two.  This is a geometric
theorem about a bosonic fluctuation; it does not yet produce a fermion.

The minimal mathematical three-mode candidate is the canonical
Spin\({}^c\) Dolbeault operator
\[
 D_T^c=\sqrt2\bigl(\bar\partial_T+\bar\partial_T^\dagger\bigr),
 \qquad
 W^+=\Omega^{0,0}(T),\quad W^-=\Omega^{0,1}(T).
\]
Hodge theory gives
\(\Ker D_T^{c,+}=H^0(\cO(2))\simeq\CC^3\),
\(\Ker D_T^{c,-}=H^1(\cO(2))=0\), and index three.  In the AP-E1 round
sphere convention \(R=1/2\), its complete spectrum is
\[
 \lambda_{n,\pm}=\pm\frac{\sqrt{n(n+3)}}{R},\qquad
 \operatorname{mult}(\lambda_{n,+})=
 \operatorname{mult}(\lambda_{n,-})=2n+3,\qquad n\ge1.
\]
Thus the first massive level is \(\pm4\), fivefold for each sign.  Only the
kernel lacks opposite chirality; every massive mode has a partner.

A decisive negative control is also proved.  The unique ordinary spin
structure has \(S^+=\cO(-1)\).  Twisting an ordinary spin Dirac operator only
by the actual tangent line \(T=\cO(2)\) leaves the Dolbeault line \(\cO(1)\)
and therefore gives two zero modes, not three.  The canonical Spin\({}^c\)
choice with \(\cO(2)\) is equivalent to ordinary spin twist \(\cO(3)\): its
extra half-canonical unit is physical input, not notation.

The companion verifier passes all exact and numerical spectral checks.
AP-E4 is therefore mathematically closed at the canonical-operator level,
but the physical origin remains open: a bosonic tangent fluctuation does not
derive a target-space fermion, a product compactification or moduli-space SQM
has not been specified, higher-dimensional anomalies are not cancelled, and
the three automorphism vector fields are not yet chiral Route-E families.
Physics promotion remains false.
\end{abstract}

\tableofcontents

\section{Question, conventions, and exact status}

AP-E4 must keep four statements separate:
\begin{enumerate}[label=(Q\arabic*)]
\item Is the linearized physical field tangent-valued?
\item Which Clifford module and Dirac operator act on it?
\item What are the chirality, complete spectrum, and first gap?
\item Under what additional microscopic conditions may the index be read as
  a four-dimensional family count?
\end{enumerate}
The first three have exact mathematical answers once the operator choice is
declared.  The fourth is a physical UV problem and is not solved by an index.

We retain the AP-E1 Fubini--Study convention
\begin{equation}
 \dd s^2=\frac{\dd w\,\dd\bar w}{(1+|w|^2)^2}
 =R^2(\dd\theta^2+\sin^2\theta\,\dd\phi^2),
 \qquad R=\frac12.
 \label{eq:metric}
\end{equation}
Consequently
\begin{equation}
 K_G=\frac1{R^2}=4,
 \qquad \operatorname{Scal}=\frac2{R^2}=8.
 \label{eq:curvatures}
\end{equation}
Every eigenvalue below refers to the canonical geometric normalization
\(D^c=\sqrt2(\bar\partial+\bar\partial^\dagger)\) and to
Eq.~\eqref{eq:metric}.

\section{First-principles derivation of the tangent bundle}

\subsection{Horizontal projection of a projective doublet}

Let a sigma-model field be represented by
\begin{equation}
 z(x)\in\CC^2,\qquad z^\dagger z=1,\qquad
 z\sim e^{\ii\alpha(x)}z.
 \label{eq:projective-field}
\end{equation}
An allowed variation obeys
\(\operatorname{Re}(z^\dagger\delta z)=0\).  Its vertical gauge component is
proportional to \(\ii z\), while the gauge-horizontal physical variation is
\begin{equation}
 \eta=Q_z\delta z,
 \qquad Q_z=1-zz^\dagger,
 \qquad z^\dagger\eta=0.
 \label{eq:horizontal-variation}
\end{equation}
Under \(z\mapsto e^{\ii\alpha}z\), one has
\(\eta\mapsto e^{\ii\alpha}\eta\).  Intrinsically,
\begin{equation}
 T_{[z]}\CP^1\simeq\operatorname{Hom}(\CC z,\CC^2/\CC z),
 \label{eq:tangent-hom}
\end{equation}
and \(\eta\) represents exactly this homomorphism.  Thus a fluctuation about
\(\varphi_0:X\to\CP^1\) lies in
\begin{equation}
 \eta\in\Gamma(X,\varphi_0^*T\CP^1).
 \label{eq:pullback-tangent}
\end{equation}
Notice the pullback: this does not yet make \(\eta\) a spinor living on target
\(\CP^1\).

\subsection{Two-chart proof that \texorpdfstring{\(T^{1,0}\CP^1=\cO(2)\)}{T CP1=O(2)}}

On \(U_N\), set
\begin{equation}
 z_N(w)=\frac{(1,w)^T}{\sqrt{1+|w|^2}}.
 \label{eq:north-frame}
\end{equation}
A direct differentiation and projection gives
\begin{equation}
 Q_N\delta z_N
 =\frac{(-\bar w,1)^T}{(1+|w|^2)^{3/2}}\,\delta w.
 \label{eq:north-horizontal}
\end{equation}
On \(U_S\), use \(v=1/w\) and
\begin{equation}
 z_S(v)=\frac{(v,1)^T}{\sqrt{1+|v|^2}},
 \qquad
 Q_S\delta z_S
 =\frac{(1,-\bar v)^T}{(1+|v|^2)^{3/2}}\,\delta v.
 \label{eq:south-horizontal}
\end{equation}
On the overlap,
\begin{equation}
 \delta v=\frac{\dd v}{\dd w}\delta w=-w^{-2}\delta w.
 \label{eq:tangent-transition}
\end{equation}
The holomorphic tangent coefficient therefore has degree two.  Hence
\begin{equation}
 \boxed{T^{1,0}\CP^1\simeq\cO(2).}
 \label{eq:tangent-o2}
\end{equation}

For the positive AP-E1 orientation, a compatible Chern connection is
\begin{align}
 a_N^{(2)}&=(1-\cos\theta)\dd\phi,
 &a_S^{(2)}&=-(1+\cos\theta)\dd\phi,\label{eq:tangent-potentials}\\
 a_N^{(2)}-a_S^{(2)}&=2\dd\phi,
 &F_T&=\sin\theta\,\dd\theta\wedge\dd\phi.
 \label{eq:tangent-curvature}
\end{align}
Thus
\begin{equation}
 \frac1{2\pi}\int_{S^2}F_T=2.
 \label{eq:tangent-chern}
\end{equation}
The transition proof fixes the algebraic bundle; the flux fixes its oriented
connection without relying on ambiguous names for ket versus dual lines.

\subsection{The bosonic Hessian is not a Dirac operator}

For
\begin{equation}
 S_B=\frac{f^2}{2}\int_X|\dd\varphi|^2,
 \label{eq:bosonic-action}
\end{equation}
the covariant exponential-map expansion gives
\begin{equation}
 S_B^{(2)}=\frac{f^2}{2}\int_X
 \left[
 \langle D_\mu\eta,D^\mu\eta\rangle
 -\langle R(\eta,\partial_\mu\varphi_0)
 \partial^\mu\varphi_0,\eta\rangle
 \right].
 \label{eq:jacobi-hessian}
\end{equation}
The Jacobi operator is second order,
\begin{equation}
 J\eta=-D_\mu D^\mu\eta
 -R(\eta,\partial_\mu\varphi_0)\partial^\mu\varphi_0
 \label{eq:jacobi-operator}
\end{equation}
in the displayed Euclidean sign convention.  For constant \(\varphi_0\), the
pullback bundle is trivial on \(X\), the curvature term vanishes, and one gets
only \(-\Box_X\).  If \(\varphi_0\) has degree \(d\) on a two-cycle, then
\(c_1(\varphi_0^*T)=2d\), not unconditionally two.

\begin{gate}[Statistics and base manifold]
The field \(\eta\) is commuting and carries no Clifford index.  A fermion
requires a separate spin or Spin\({}^c\) module.  In an ordinary supersymmetric
sigma model, \(\psi\in\Gamma(S_X\otimes\varphi^*T)\) is acted on by a Dirac
operator on the base \(X\), not automatically by a Dirac operator whose
coordinate is target \(\CP^1\).  A target spectrum requires an explicit
moduli-space quantum mechanics or a product compactification.
\end{gate}

\section{The canonical Spin\texorpdfstring{\({}^c\)}{c} tangent operator}

\subsection{Operator, chirality, and global bundle}

On a K\"ahler curve, the canonical Spin\({}^c\) module is
\cite{LawsonMichelsohn1989,Wells2008}
\begin{equation}
 W^+=\Omega^{0,0},\qquad W^-=\Omega^{0,1},
 \qquad \Gamma|_{W^\pm}=\pm1.
 \label{eq:canonical-module}
\end{equation}
Twist it by
\(E=T^{1,0}\CP^1=\cO(2)\).  The selected AP-E4 operator is
\begin{equation}
 \boxed{
 D_T^c=\sqrt2(\bar\partial_T+\bar\partial_T^\dagger):
 \Omega^{0,0}(T)\oplus\Omega^{0,1}(T)
 \longrightarrow
 \Omega^{0,0}(T)\oplus\Omega^{0,1}(T).}
 \label{eq:canonical-dirac}
\end{equation}
It is self-adjoint and odd,
\begin{equation}
 \{D_T^c,\Gamma\}=0,
 \qquad (D_T^c)^2=2\Delta_{\bar\partial_T}.
 \label{eq:dirac-square}
\end{equation}

The canonical Spin\({}^c\) determinant before twisting is
\(K^{-1}=\cO(2)\).  Twisting the module by \(E=\cO(q)\) changes the effective
determinant to
\begin{equation}
 \det W_E=K^{-1}\otimes E^2=\cO(2q+2).
 \label{eq:spinc-determinant}
\end{equation}
For \(q=2\), this is \(\cO(6)\).  This topological datum must eventually be
part of the microscopic theory.

\subsection{Local ordinary-spin representative and the hidden shift}

The unique spin structure on \(S^2\) has
\begin{equation}
 S^+=K^{1/2}=\cO(-1),\qquad S^-=K^{-1/2}=\cO(1).
 \label{eq:spin-bundles}
\end{equation}
An ordinary spin Dirac operator twisted by \(\cO(p)\) may be written on the
north patch as
\begin{equation}
 D_p^{\rm spin}=-\frac{\ii}{R}\left[
 \sigma_1\left(\partial_\theta+\frac12\cot\theta\right)
 +\frac{\sigma_2}{\sin\theta}
 \left(\partial_\phi-\ii a_\phi^{(p)}\right)
 \right],
 \label{eq:local-spin-dirac}
\end{equation}
where
\begin{equation}
 a_N^{(p)}=\frac p2(1-\cos\theta)\dd\phi,
 \qquad
 a_S^{(p)}=-\frac p2(1+\cos\theta)\dd\phi.
 \label{eq:op-connection}
\end{equation}
Its positive component is the Dolbeault line
\begin{equation}
 S^+\otimes\cO(p)=\cO(p-1).
 \label{eq:ordinary-effective-line}
\end{equation}
Consequently
\begin{equation}
 D^c_{\cO(q)}\simeq D^{\rm spin}_{\cO(q+1)}.
 \label{eq:spin-spinc-relation}
\end{equation}
The canonical tangent choice \(q=2\) is equivalent to ordinary twist
\(p=3\), not \(p=2\).  The unit shift is the half-canonical factor
\(K^{-1/2}=\cO(1)\).

\begin{gate}[Ordinary tangent negative control]
If a microscopic fermion is an ordinary spinor twisted only by the actual
tangent line \(T=\cO(2)\), then \(p=2\) and its positive Dolbeault line is
\(\cO(1)\).  It has two zero modes.  Calling this operator
``Spin\({}^c\) on \(T\)'' does not change the count.  The physical problem is
to derive the canonical Spin\({}^c\) module or the extra \(\cO(1)\), not to
rename it.
\end{gate}

\section{Index, chirality, and exact zero modes}

Hodge theory gives \cite{Wells2008}
\begin{equation}
 \Ker D_T^{c,+}\simeq H^0(\CP^1,\cO(2)),
 \qquad
 \Ker D_T^{c,-}\simeq H^1(\CP^1,\cO(2)).
 \label{eq:hodge-kernels}
\end{equation}
For \(\cO(q)\) on \(\CP^1\),
\begin{equation}
 h^0(\cO(q))=\max\{q+1,0\},
 \qquad
 h^1(\cO(q))=h^0(\cO(-q-2)).
 \label{eq:line-cohomology}
\end{equation}
The second identity is Serre duality.  At \(q=2\),
\begin{equation}
 h^0(\cO(2))=3,
 \qquad h^1(\cO(2))=h^0(\cO(-4))=0.
 \label{eq:q2-cohomology}
\end{equation}
Equivalently, by the index theorem \cite{AtiyahSinger1968}, with
\(x=c_1(\cO(1))\) and \(\int_{\CP^1}x=1\),
\begin{align}
 \ind D_T^c
 &=\int_{\CP^1}\operatorname{ch}(\cO(2))
 \operatorname{Td}(T\CP^1)\notag\\
 &=\int_{\CP^1}(1+2x)(1+x)=3.
 \label{eq:index-theorem}
\end{align}
Thus
\begin{equation}
 \boxed{\dim\Ker D_T^{c,+}=3,\quad
 \dim\Ker D_T^{c,-}=0,\quad \ind D_T^c=3.}
 \label{eq:three-zero-modes}
\end{equation}
The basis is the familiar
\(\{z_0^2,z_0z_1,z_1^2\}\), and
\(H^0(T\CP^1)\simeq\mathfrak{sl}_2(\CC)\) as a complex representation.

\begin{remark}[Orientation reversal is shifted]
For the canonical Spin\({}^c\) operator, the index of \(\cO(q)\) is \(q+1\).
Three negative-chirality zero modes require \(q=-4\), not \(q=-2\).  This
asymmetry is another manifestation of the built-in canonical determinant.
An ordinary spin monopole spectrum has the more familiar sign-symmetric
charge count after its half-canonical shift is included explicitly.
\end{remark}

\section{Complete spectrum and spectral gap}

\subsection{Monopole harmonics and the Weitzenb\"ock shift}

For a canonical Dolbeault twist \(E=\cO(q)\), \(q\ge0\), the relevant
\(SU(2)\) multiplet at excitation number \(n\ge0\) has
\begin{equation}
 j=\frac q2+n,
 \qquad \dim V_j=2j+1=q+2n+1.
 \label{eq:monopole-j}
\end{equation}
The connection Laplacian has eigenvalue
\begin{equation}
 \nabla^*\nabla
 =\frac1{R^2}\left[j(j+1)-\left(\frac q2\right)^2\right].
 \label{eq:bochner-spectrum}
\end{equation}
On positive chirality, the K\"ahler--Weitzenb\"ock identity subtracts the
curvature term \(q/(2R^2)\):
\begin{equation}
 (D^c_{\cO(q)})^2
 =2\bar\partial_q^\dagger\bar\partial_q
 =\nabla^*\nabla-\frac{q}{2R^2}.
 \label{eq:weitzenbock}
\end{equation}
Substitution of Eq.~\eqref{eq:monopole-j} gives the exact identity
\begin{align}
 j(j+1)-\frac{q^2}{4}-\frac q2
 &=n(n+q+1).
 \label{eq:casimir-identity}
\end{align}
At \(n=0\) this vanishes and the multiplicity is \(q+1\), reproducing the
index result.

For \(n\ge1\), supersymmetry/Hodge decomposition pairs the positive and
negative form sectors.  Since \(D\) is odd and self-adjoint, every positive
eigenvalue has a negative partner.  The full spectrum is therefore
\begin{equation}
 \boxed{
 \lambda_{n,\pm}=\pm\frac{\sqrt{n(n+q+1)}}{R},\qquad
 \operatorname{mult}(\lambda_{n,\pm})=q+2n+1,\qquad n\ge1.}
 \label{eq:full-q-spectrum}
\end{equation}
This agrees with the supersymmetric projective-space monopole spectrum after
converting operator and radius normalizations \cite{BykovSmilga2023APE4}.

\subsection{The tangent result at \texorpdfstring{\(q=2,R=1/2\)}{q=2,R=1/2}}

For AP-E4,
\begin{equation}
 \lambda_{n,\pm}=\pm2\sqrt{n(n+3)},
 \qquad \operatorname{mult}(\lambda_{n,\pm})=2n+3.
 \label{eq:tangent-spectrum}
\end{equation}
The first levels are
\begin{center}
\begin{tabular}{@{}cccc@{}}
\toprule
\(n\) & \(|\lambda_n|\) & multiplicity & interpretation\\
\midrule
0 & 0 & \(3\) in \(W^+\), \(0\) in \(W^-\) & chiral kernel; no
\(\pm\lambda\) split\\
1 & 4 & 5 & first massive partner level\\
2 & \(2\sqrt{10}\) & 7 & massive partner level\\
3 & \(6\sqrt2\simeq8.48528\) & 9 & massive partner level\\
\bottomrule
\end{tabular}
\end{center}
Thus
\begin{equation}
 \boxed{\Delta_D=4,\qquad \Delta_{D^2}=16.}
 \label{eq:canonical-gap}
\end{equation}
If the microscopic kinetic operator carries a velocity or normalization
\(v_D\), all masses are multiplied by \(v_D\); the dimensionless geometric
gap is the value above.

\subsection{Ordinary spin twisted by the tangent line}

For ordinary \(p=2\), the effective Dolbeault degree is \(q=p-1=1\).  Hence
\begin{equation}
 \dim\Ker D_{p=2}^{\rm spin,+}=2,\qquad
 \lambda_{n,\pm}^{\rm spin}
 =\pm\frac{\sqrt{n(n+p)}}{R},\qquad
 \operatorname{mult}=p+2n.
 \label{eq:ordinary-spectrum}
\end{equation}
At \(R=1/2\),
\begin{equation}
 \Delta_{p=2}^{\rm spin}=2\sqrt3\simeq3.464101615,
 \qquad \operatorname{mult}_{n=1}=4
 \label{eq:ordinary-gap}
\end{equation}
per sign.  This different kernel, gap, and degeneracy make the hidden
Spin\({}^c\) choice experimentally and numerically falsifiable.

\section{Partner spectrum and possible four-dimensional chirality}

\subsection{What ``no opposite chirality'' actually means}

Equation~\eqref{eq:three-zero-modes} proves that the \emph{kernel} has no
negative-chirality state.  It does not remove negative chirality from the
theory.  For every \(n\ge1\), choose normalized chiral components
\(\psi_n^{(0)}\in W^+\) and \(\psi_n^{(1)}\in W^-\) so that
\begin{equation}
 D_T^c\psi_n^{(0)}=|\lambda_n|\psi_n^{(1)},\qquad
 D_T^c\psi_n^{(1)}=|\lambda_n|\psi_n^{(0)}.
 \label{eq:chiral-component-pair}
\end{equation}
The operator then produces the normalized eigenstates
\begin{equation}
 \psi_{n,+}=\frac{\psi_n^{(0)}+\psi_n^{(1)}}{\sqrt2},\qquad
 \psi_{n,-}=\frac{\psi_n^{(0)}-\psi_n^{(1)}}{\sqrt2},
 \label{eq:massive-pairs}
\end{equation}
with eigenvalues \(\pm|\lambda_n|\).  The first massive level contains five
such pairs.  Any portal with characteristic scale comparable to \(4v_D\) will
mix these states into the purported three-mode sector.

\subsection{A conditional product compactification}

Only after declaring a product geometry \(M^{3,1}\times\CP^1\) may one write
schematically
\begin{equation}
 \slashed D_6=\slashed D_4\otimes1
 +\gamma_5\otimes D_T^c,
 \qquad
 \Gamma_7=\gamma_5\otimes\Gamma_2,
 \label{eq:product-dirac}
\end{equation}
up to a conventional overall sign.  A single six-dimensional Weyl condition
then ties internal to four-dimensional chirality: the three internal
positive zero modes become three same-chirality four-dimensional massless
modes.  Each nonzero internal pair becomes a massive four-dimensional Dirac
pair.  A six-dimensional \emph{Dirac} field contains both six-dimensional
chiralities and generally doubles the four-dimensional zero sector.

\begin{gate}[Higher-dimensional anomalies]
A single six-dimensional Weyl fermion is generally gauge- and
gravitationally anomalous.  A physical compactification must cancel its full
eight-form anomaly polynomial, or exhibit a valid Green--Schwarz
factorization, and must quantize the Spin\({}^c\) determinant flux.  After
reduction, the representation-weighted four-dimensional anomaly is multiplied
by the index three.  These are spectrum-level consistency conditions, not
optional phenomenology \cite{AlvarezGaumeWitten1984}.
\end{gate}

\subsection{The automorphism-versus-family gate}

The vector space
\begin{equation}
 H^0(\CP^1,T\CP^1)\simeq\mathfrak{sl}_2(\CC)
 \label{eq:automorphism-space}
\end{equation}
consists of holomorphic infinitesimal automorphisms of \(\CP^1\).  Depending
on the microscopic construction, these may be gauge redundancies, collective
coordinates, or physical matter modes.  Dimension three alone cannot decide.
One must derive the kinetic norm, gauge quotient, couplings, and anomaly
representation of each mode before calling them three families.

\section{Deterministic analytical and numerical audit}

The companion script
\texttt{route\_f/code/verify\_ap\_e4\_tangent\_dirac\_spectrum.py}
is deliberately independent of a fitted spectrum.  It performs:
\begin{enumerate}
\item exact line-bundle cohomology and Serre-duality checks for \(q=2\),
  ordinary \(p=2\), and neighbouring \(q=1,3\) controls;
\item exact rational verification of Eq.~\eqref{eq:casimir-identity} and its
  multiplicity through ten levels;
\item numerical construction of each finite \(SU(2)\) spin-\(j\) matrix,
  diagonalization of the shifted Casimir, and comparison with
  \(n(n+q+1)\);
\item a truncated chiral block operator with three positive zero bases and
  five massive levels, testing Hermiticity,
  \(\{D,\Gamma\}=0\), exact \(\pm\lambda\) pairing, kernel chirality, and
  the gap four;
\item curvature quadrature for \(c_1(T)=2\), horizontal-projection residuals,
  and radius-rescaling checks;
\item negative controls for nonintegral monopole charge, ordinary-spin
  misidentification, a chirality-breaking mass, and the false assertion that
  massive partners are absent.
\end{enumerate}
The spectral block regression checks bookkeeping, while the independent
Casimir diagonalization checks the representation-theoretic eigenvalues.
Neither is presented as a substitute for a microscopic origin.

\section{Claim ledger and next steps}

\begin{longtable}{@{}L{0.43\textwidth}L{0.20\textwidth}L{0.28\textwidth}@{}}
\toprule
Claim & Status & Boundary\\
\midrule
\endfirsthead
\toprule
Claim & Status & Boundary\\
\midrule
\endhead
Horizontal projective fluctuation is tangent-valued & \status{proved} & as a
section of \(\varphi_0^*T\CP^1\)\\
\(T^{1,0}\CP^1=\cO(2)\), \(c_1=2\) & \status{proved} & chart and flux proofs\\
Canonical Spin\({}^c\) tangent operator has three positive zero modes &
\status{proved} & declared operator choice\\
Complete canonical spectrum and gap four & \status{proved} & AP-E1
\(R=1/2\) convention\\
Every massive level has paired \(\pm D_T^c\) eigenvalues and equal
\(W^\pm\) multiplicities & \status{proved} & only the kernel has a chiral
imbalance\\
Ordinary spin twist by \(T\) has two zero modes & \status{proved} & decisive
negative control\\
Physical fermion is canonical Spin\({}^c\) and tangent-valued &
\status{open} & requires moduli SQM or compactification\\
Six-dimensional and four-dimensional anomalies cancel & \status{open} & no
complete field content or \(I_8\) factorization\\
Three zero modes are three Route-E families & \status{open} & automorphism,
gauge, and portal gates unresolved\\
\bottomrule
\end{longtable}

Three physically distinct continuations should be retained:
\begin{enumerate}
\item \textbf{Moduli-space/SQM route.}  Derive a soliton moduli space
  \(\CP^1\), its tangent fermions, and \(N=2\) supercharge
  \(Q=\sqrt2\bar\partial_T\); then compute whether these quantum states carry
  the required four-dimensional gauge representation.
\item \textbf{Product-compactification route.}  Specify a six-dimensional
  anomaly-free theory on \(M^{3,1}\times\CP^1\), derive the Spin\({}^c\)
  determinant line and \(E=T\), and calculate the complete KK/threshold
  spectrum.
\item \textbf{Ordinary-spin control route.}  Accept the natural ordinary
  tangent twist and its two modes unless an additional \(\cO(1)\) gauge flux
  is independently derived.  This is the conservative fallback.
\end{enumerate}

The fail-closed result is
\begin{equation}
 \boxed{\begin{gathered}
 \texttt{tangent\_bundle\_derived=true},\\
 \texttt{canonical\_spinc\_mathematics\_complete=true},\\
 \texttt{canonical\_spinc\_three\_zero\_modes=true},\\
 \texttt{partner\_spectrum\_and\_gap\_complete=true},\\
 \texttt{physical\_fermionic\_tangent\_mode\_derived=false},\\
 \texttt{six\_dimensional\_anomaly\_cancelled=false},\\
 \texttt{route\_e\_chiral\_family\_identification=false},\\
 \texttt{ap\_e4\_physics\_closed=false},\\
 \texttt{physics\_promotion\_allowed=false}.
 \end{gathered}}
 \label{eq:final-status}
\end{equation}

\bibliographystyle{unsrt}
\bibliography{route_f/tex/ap_e4_tangent_dirac_spectrum}

\end{document}
